\introduction

El panorama de la Educación Superior en el siglo XXI está inmerso en una 
dinámica de cambio sin precedentes por rápidos avances tecnológicos y cambios sociales.
La aceleración tecnológica y las fluctuantes demandas del mercado laboral exigen que las instituciones
académicas migren de currículos rígidos hacia planes de estudio dinámicos y
flexibles \citep{unesco2022}. Esta necesidad de flexibilidad introduce un desafío
crítico: la gestión de la variabilidad curricular. 

La innovación educativa y la gestión curricular 
se han convertido en elementos clave para transformar el paradigma educativo y 
preparar a los estudiantes para un futuro cada vez más digitalizado y globalizado. Ambas son 
aspectos fundamentales para el desarrollo de sistemas educativos actuales, 
sin embargo, en el presente, resalta la necesidad de una gestión curricular
que promueva la colaboración entre docentes, la adaptación a las necesidades del alumnado
y la actualización de los contenidos, métodos pedagógicos, planes de estudios, 
planificación de asignaturas e incluso, carreras en general.

La innovación educativa se refiere a la introducción de nuevas ideas, enfoques,
metodologías y tecnologías en el ámbito educativo; busca mejorar la calidad de la enseñanza
y el aprendizaje, fomentando la creatividad, el pensamiento crítico y la resolución de
problemas en los estudiantes. La innovación educativa también promueve la adaptación de
los procesos educativos a las necesidades cambiantes de la sociedad y el entorno laboral \citep{gomez2023}.
Por otro lado, la gestión curricular implica el diseño, desarrollo e implementación de planes
de estudio y programas educativos. Esta última, se ocupa de
organizar y estructurar los contenidos, las habilidades y las competencias que se enseñan
en las instituciones educativas; además, implica la evaluación y la mejora contínua de los
programas curriculares para asegurar su relevancia y efectividad \citep{gomez2023}. 

La innovación educativa y la gestión curricular están estrechamente interrelacionadas.
La introducción de enfoques innovadores en la enseñanza y el aprendizaje requiere una
gestión curricular flexible y adaptable, esto implica revisar y actualizar constantemente los
planes de estudio para integrar nuevas metodologías, tecnologías y recursos educativos.
Igualmente, la gestión curricular efectiva es fundamental para garantizar que los cambios e
innovaciones en el ámbito educativo se implementen de manera coherente y sistemática. La
integración de las \ac{TIC} en el aula permite
ampliar el acceso a la información, fomentar la colaboración entre estudiantes, personalizar
el aprendizaje y desarrollar habilidades digitales clave, de la misma manera, la gestión
curricular puede incluir la planificación y el uso estratégico de la tecnología en los
programas educativos para maximizar su impacto y beneficios. \citep{gomez2023}

La búsqueda de esta flexibilidad ha llevado a la exploración de paradigmas de ingeniería de software aplicados
al ámbito educativo. El concepto de \ac{SPL}
ofrece un enfoque estructurado y probado para gestionar familias de productos
a partir de un conjunto común de activos \citep{apel2013feature}. Cuando este
paradigma se aplica al dominio educativo, surge el concepto de "Líneas de
Productos Educativos", donde un tronco curricular común puede derivar en
múltiples planes de estudio especializados mediante la gestión explícita de su
variabilidad. La investigación en este campo es incipiente pero prometedora;
por ejemplo, estudios en el ecosistema de software educativo \textit{SOLAR} han
demostrado que la modelación de la variabilidad, incluso en componentes
dinámicos como los foros de discusión, es fundamental para crear sistemas
adaptativos que respondan a contextos educativos diversos \citep{guzman2022variabilidad}.

Actualmente, la definición y el mantenimiento de la variabilidad curricular se gestionan 
principalmente mediante documentos estáticos, hojas de cálculo y procesos manuales, un enfoque que resulta 
costoso, lento y altamente propenso a errores. Esta dependencia de herramientas fragmentadas y no integradas 
dificulta la actualización oportuna de la información, genera inconsistencias entre versiones y 
complica la trazabilidad de los cambios realizados. Como consecuencia, pueden diseñarse rutas formativas 
incoherentes o inválidas, afectando la progresión del estudiante y la calidad del proceso educativo. Este tipo de 
gestión incrementa la carga de trabajo administrativo y docente, limita la visibilidad global 
del currículo, provoca la divulgación de información desactualizada y compromete la coherencia institucional. 
Todo ello tiene un impacto directo en la satisfacción del estudiante, en los procesos de acreditación y en la 
capacidad de la institución para adaptarse rápidamente a nuevas demandas educativas o del mercado laboral. 
En conjunto, estos efectos negativos evidencian que los métodos tradicionales basados en documentación 
manual ya no son sostenibles, pues ponen en riesgo la calidad, pertinencia y eficiencia del modelo formativo.

La proliferación de modalidades educativas (presenciales, híbridas, en línea,
\ac{MOOC}, \textit{micro-learning}) y enfoques pedagógicos (aprendizaje basado en
proyectos, \textit{flipped classroom}, gamificación) ha creado un panorama de extrema
variabilidad instruccional \citep{bates2019teaching}. Actualmente, esta complejidad se
gestiona predominantemente mediante soluciones \textit{ad-hoc} que resultan en
implementaciones rígidas y de alto costo de mantenimiento \citep{sanchez2020metodologia}.
La educación se enfrenta a la difícil actividad de crear experiencias de
aprendizaje verdaderamente personalizadas para cada estudiante o contexto
institucional sin un esfuerzo manual enorme. Irónicamente, la misma tecnología que ha posibilitado esta 
explosión de modalidades y enfoques pedagógicos carece, en la actualidad, de los mecanismos sistemáticos para 
gestionar la complejidad que ella misma ha generado, pues es más notable la ausencia de un enfoque 
sistemático para gestionar la variabilidad educativa de las instituciones de Educación Superior.

Cada nuevo curso o recurso educativo se desarrolla practicamente desde cero, sin reutilizar componentes
comunes de manera sistemática. Es innegable que existe a día de hoy una batalla contrarreloj,
dado que la velocidad de la innovación tecnológica supera con creces la capacidad de reacción de los 
sistemas educativos tradicionales, volviendo potencialmente obsoleta cualquier renovación curricular por parte 
de las instituciones académicas, incluso antes de su implementación o, en el peor de los casos,
incluso en el momento de su lanzamiento. Esta forma de operar no solo resulta insostenible a largo plazo, 
sino que desaprovecha la oportunidad de construir conocimiento educativo de manera incremental, 
donde cada nuevo desarrollo pueda partir de una base sólida y previamente validada. Queda así de manifiesto que 
el paradigma artesanal de diseño curricular, curso por curso, ha llegado a su límite, exigiendo un 
cambio estructural hacia modelos que garanticen tanto la agilidad como la coherencia pedagógica en la 
formación de profesionales.

De este análisis, se desprende la siguiente interrogante que define el
\textbf{problema de investigación:}
¿Cómo implementar un servicio \textit{backend} que proporcione la funcionalidad completa
para la gestión, validación y configuración de \ac{FM} en el contexto de Líneas de Productos Educativos?
Tomando como \textbf{objeto de estudio:} 
\ac{SPLE} aplicada al dominio educativo.
Enmarcado en el \textbf{campo de acción:}
servicios \textit{backend} especializados para la gestión de la variabilidad en Líneas de Productos Educativos.
Con el fin de solucionar la problemática planteada, se define como \textbf{objetivo general:}
desarrollar un servicio \textit{backend} para la gestión de la variabilidad
en Líneas de Productos Educativos.

Para dar cumplimiento al objetivo planteado se proponen las siguientes \textbf{tareas investigativas:}
\begin{enumerate}
    \item \textbf{Estudio} exhaustivo del estado del arte en Ingeniería de
    Líneas de Productos de Software, \ac{FM} y su
    aplicación al dominio educativo.
    \item \textbf{Análisis} de diferentes soluciones homólogas para la identificación de características
    generales y tendencias en este tipo de aplicaciones.
    \item \textbf{Diseño} de la arquitectura del servicio \textit{backend}, definiendo los
    componentes principales  y el modelo de datos para
    representar \ac{FM} en el contexto educativo.
    \item \textbf{Selección} de las herramientas, tecnologías y metodología a emplear, para garantizar
    una excelente y eficiente experiencia de desarrollo y alineado con los objetivos del proyecto.
    \item \textbf{Identificación} de los requisitos funcionales y no funcionales necesarios para un
    sistema capaz de gestionar la variabilidad de Líneas de Productos Educativos.
    \item \textbf{Implementación} del servicio \textit{backend} para materializar el sistema que dará respuesta
    a la problemática planteada, asegurando la creación, análisis, validación y configuración de \ac{FM}.
    \item \textbf{Validación} del servicio desarrollado mediante casos de uso representativos
    en el ámbito educativo, evaluando su funcionalidad, rendimiento, utilidad y 
    efectividad para capturar la variabilidad en la educación.
    \item \textbf{Documentación} de los resultados obtenidos y elaboración de recomendaciones para
    futuras investigaciones en el área.
\end{enumerate}

Para darle solución a los objetivos trazados en las tareas de investigación se emplearon los
siguientes \textbf{métodos científicos:}

\begin{itemize}
    \item \textbf{Métodos teóricos:} 
        \begin{itemize}
            \item \textbf{Modelación:} se utilizó para la elaboración de los diagramas del lenguaje de
            modelado y en la representación de las características de la solución.
            \item \textbf{Histórico-Lógico:} Se empleó para comprender la evolución de las técnicas de
            gestión de variabilidad en \ac{SPL}, desde sus orígenes en ingeniería de software
            hasta su aplicación potencial en dominios educativos, analizando su desarrollo
            histórico y lógica interna.
            \item \textbf{Analítico-sintético:} Utilizado en el examen crítico de la bibliografía
            especializada sobre \ac{SPL}, \ac{FM} y sistemas educativos, permitiendo
            sintetizar un marco conceptual adaptado al dominio educativo.
        \end{itemize}
    \item \textbf{Métodos empíricos:}
        \begin{itemize}
            \item \textbf{Observación:} Mediante este método se analizaron herramientas
            existentes de gestión \ac{SPL} (como \textit{pure::variants}, \textit{FeatureIDE}) para
            identificar limitaciones técnicas y requisitos necesarios para el contexto
            educativo específico.
            \item \textbf{Entrevista:} Se entrevista especialistas responsables de diseños
            curriculares que compran aspectos técnicos para identificar los desafíos
            existentes en la flexibilidad y variabilidad de los currículos.
            \item \textbf{Experimentación:} Mediante la implementación de prototipos
            funcionales para validar diferentes enfoques arquitectónicos y algoritmos
            de validación.
        \end{itemize}
\end{itemize}
